% ===============================================================================
% CHAPTER 2: DATA CLEANING AND PREPROCESSING
% ===============================================================================

\subsection{Data Cleaning and Preprocessing}

This chapter presents the systematic approach used to load, assess, and clean the 3D printer dataset. The data preprocessing pipeline ensures data quality and prepares the dataset for subsequent statistical analysis.

\subsubsection{Data Loading}

The first step involves loading the dataset and examining its basic structure. The following R code demonstrates the data loading process:

\begin{lstlisting}[language=R, caption={Data Loading and Initial Exploration}, label={lst:data_load}]
# Load necessary libraries
library(readr)
library(dplyr)

# Load the 3D printer dataset
data <- read_csv("/home/kari/prob-stat-pj/3dprintdata/data.csv")

# Display basic information about the dataset
str(data)
head(data, 10)
colnames(data)
summary(data)
\end{lstlisting}

The dataset contains 50 observations with 12 variables, including 9 input parameters (layer\_height, wall\_thickness, infill\_density, infill\_pattern, nozzle\_temperature, bed\_temperature, print\_speed, material, fan\_speed) and 3 output quality metrics (roughness, tension\_strength, elongation). The \texttt{str()} function reveals the data types, while \texttt{summary()} provides basic descriptive statistics for each variable.

\textbf{Output:}
\begin{verbatim}
Dimensions: 50 rows x 12 columns
Variables: layer_height, wall_thickness, infill_density, infill_pattern, 
          nozzle_temperature, bed_temperature, print_speed, material, 
          fan_speed, roughness, tension_strenght, elongation
\end{verbatim}

\subsubsection{Data Quality Assessment}

Before proceeding with analysis, it is essential to assess the quality of the dataset by checking for missing values, duplicate records, and potential outliers.

\paragraph{Missing Values and Duplicates}

\begin{lstlisting}[language=R, caption={Missing Values and Duplicate Detection}, label={lst:quality_check}]
# Check for missing values
missing_values <- sapply(data, function(x) sum(is.na(x)))
print(missing_values)
cat("Total missing values:", sum(missing_values), "\n")

# Check for duplicate rows
duplicate_count <- sum(duplicated(data))
cat("Number of duplicate rows:", duplicate_count, "\n")
\end{lstlisting}

The quality assessment revealed that the dataset is complete with no missing values (0 missing values across all variables) and no duplicate records (0 duplicate rows), indicating high data quality.

\textbf{Output:}
\begin{verbatim}
Missing values: 0
Duplicate rows: 0
\end{verbatim}

\paragraph{Outlier Detection}

Outliers were detected using the Interquartile Range (IQR) method, which identifies values that fall outside the range [Q1 - 1.5×IQR, Q3 + 1.5×IQR]:

\begin{lstlisting}[language=R, caption={Outlier Detection Using IQR Method}, label={lst:outlier_detection}]
# Define continuous variables for outlier detection
continuous_vars <- c("layer_height", "nozzle_temperature", "bed_temperature", 
                    "print_speed", "infill_density", "wall_thickness", 
                    "fan_speed", "roughness", "tension_strenght", "elongation")

# Outlier detection for each continuous variable
for(var in continuous_vars) {
  if(var %in% colnames(data)) {
    Q1 <- quantile(data[[var]], 0.25, na.rm = TRUE)
    Q3 <- quantile(data[[var]], 0.75, na.rm = TRUE)
    IQR_val <- Q3 - Q1
    lower_bound <- Q1 - 1.5 * IQR_val
    upper_bound <- Q3 + 1.5 * IQR_val
    
    outliers <- data[[var]] < lower_bound | data[[var]] > upper_bound
    outlier_count <- sum(outliers, na.rm = TRUE)
    
    cat(sprintf("%s: Outliers = %d\n", var, outlier_count))
  }
}
\end{lstlisting}

The IQR method systematically evaluates each continuous variable by calculating the first quartile (Q1), third quartile (Q3), and the interquartile range. Values falling outside the calculated bounds are flagged as potential outliers for further investigation.

\subsubsection{Data Cleaning Process}

The data cleaning process involves handling any identified issues and preparing the data for analysis through appropriate data type conversions.

\paragraph{Data Type Conversions}

Categorical variables require conversion to factor type for proper statistical analysis:

\begin{lstlisting}[language=R, caption={Data Type Conversions}, label={lst:data_conversion}]
# Create a copy of original data for cleaning
data_cleaned <- data

# Convert categorical variables to factors
if("infill_pattern" %in% colnames(data_cleaned)) {
  data_cleaned$infill_pattern <- as.factor(data_cleaned$infill_pattern)
  cat("Converted infill_pattern to factor\n")
}

if("material" %in% colnames(data_cleaned)) {
  data_cleaned$material <- as.factor(data_cleaned$material)
  cat("Converted material to factor\n")
}

# Verify the conversion
str(data_cleaned)
\end{lstlisting}

The \texttt{infill\_pattern} variable contains two levels (grid and honeycomb), while the \texttt{material} variable has two levels (ABS and PLA). Converting these to factors ensures they are treated as categorical variables in subsequent statistical analyses rather than character strings.

\textbf{Output:}
\begin{verbatim}
Factor variables: infill_pattern, material
infill_pattern levels: grid (25), honeycomb (25)
material levels: abs (25), pla (25)
\end{verbatim}

\paragraph{Data Validation}

After cleaning, it is important to validate the data ranges and distributions:

\begin{lstlisting}[language=R, caption={Data Range Validation}, label={lst:data_validation}]
# Check data ranges for continuous variables
for(var in continuous_vars) {
  if(var %in% colnames(data_cleaned)) {
    cat(sprintf("%s: Min=%.3f, Max=%.3f, Range=%.3f\n", 
                var, min(data_cleaned[[var]], na.rm=TRUE), 
                max(data_cleaned[[var]], na.rm=TRUE),
                max(data_cleaned[[var]], na.rm=TRUE) - min(data_cleaned[[var]], na.rm=TRUE)))
  }
}

# Check categorical variable distributions
print(table(data_cleaned$infill_pattern))
print(table(data_cleaned$material))
\end{lstlisting}

This validation step ensures that all variables have reasonable ranges and that categorical variables have balanced distributions, which is important for the reliability of subsequent statistical analyses.

\subsubsection{Cleaned Dataset Summary}

The data cleaning process resulted in a high-quality dataset ready for analysis:

\begin{lstlisting}[language=R, caption={Final Dataset Summary}, label={lst:final_summary}]
# Compare original and cleaned datasets
cat("Original dataset:", nrow(data), "rows x", ncol(data), "columns\n")
cat("Cleaned dataset:", nrow(data_cleaned), "rows x", ncol(data_cleaned), "columns\n")
cat("Rows removed:", nrow(data) - nrow(data_cleaned), "\n")

# Final summary statistics
summary(data_cleaned)

# Save cleaned dataset
write_csv(data_cleaned, "/home/kari/prob-stat-pj/3dprintdata/data_cleaned.csv")
\end{lstlisting}

\paragraph{Key Findings from Data Cleaning}

The data cleaning process revealed several important characteristics:

\begin{itemize}
    \item \textbf{Data Completeness}: The dataset is complete with no missing values across all 50 observations and 12 variables
    \item \textbf{Data Integrity}: No duplicate records were found, ensuring each observation represents a unique experimental condition
    \item \textbf{Variable Types}: Successfully converted 2 categorical variables (infill\_pattern and material) to factors with balanced distributions (25 observations each for grid/honeycomb patterns and ABS/PLA materials)
    \item \textbf{Data Quality}: All continuous variables show reasonable ranges consistent with typical 3D printing parameters and quality metrics
    \item \textbf{Final Dataset}: The cleaned dataset maintains the original 50 × 12 structure, indicating no data loss during the cleaning process
\end{itemize}

\textbf{Summary Statistics of Cleaned Dataset:}
\begin{verbatim}
  layer_height   wall_thickness  infill_density   infill_pattern
 Min.   :0.020   Min.   : 1.00   Min.   :10.0   grid     :25    
 1st Qu.:0.060   1st Qu.: 3.00   1st Qu.:40.0   honeycomb:25    
 Median :0.100   Median : 5.00   Median :50.0                   
 Mean   :0.106   Mean   : 5.22   Mean   :53.4                   
 3rd Qu.:0.150   3rd Qu.: 7.00   3rd Qu.:80.0                   
 Max.   :0.200   Max.   :10.00   Max.   :90.0                   

 nozzle_temperature bed_temperature  print_speed  material   fan_speed  
 Min.   :200.0      Min.   :60      Min.   : 40   abs:25   Min.   :  0  
 1st Qu.:210.0      1st Qu.:65      1st Qu.: 40   pla:25   1st Qu.: 25  
 Median :220.0      Median :70      Median : 60            Median : 50  
 Mean   :221.5      Mean   :70      Mean   : 64            Mean   : 50  
 3rd Qu.:230.0      3rd Qu.:75      3rd Qu.: 60            3rd Qu.: 75  
 Max.   :250.0      Max.   :80      Max.   :120            Max.   :100  

   roughness     tension_strenght   elongation   
 Min.   : 21.0   Min.   : 4.00    Min.   :0.400  
 1st Qu.: 92.0   1st Qu.:12.00    1st Qu.:1.100  
 Median :165.5   Median :19.00    Median :1.550  
 Mean   :170.6   Mean   :20.08    Mean   :1.672  
 3rd Qu.:239.2   3rd Qu.:27.00    3rd Qu.:2.175  
 Max.   :368.0   Max.   :37.00    Max.   :3.300  
\end{verbatim}

The cleaned dataset provides a solid foundation for descriptive statistics and subsequent multiple linear regression analysis, with all variables properly formatted and validated for statistical modeling.